\documentclass[french,a4paper]{article}

\usepackage[utf8]{inputenc}
\usepackage[T1]{fontenc}
\usepackage{lmodern}
\usepackage{geometry}
\usepackage{graphicx}
\usepackage{babel}
\usepackage{amsmath}
\usepackage{amsthm}
\usepackage{amssymb}
\usepackage{hyperref}
\usepackage{stmaryrd}
\usepackage{enumitem}
\usepackage{tcolorbox}
\usepackage{dsfont}
\usepackage{multirow}
\usepackage{etoc}
\usepackage{titlesec}
\usepackage{esint}
\usepackage{cancel}
\usepackage{mathdots}

\setlist[itemize]{label=$\bullet$}


\renewcommand{\P}{\mathbb{P}}
\newcommand{\N}{\mathbb{N}}
\newcommand{\Z}{\mathbb{Z}}
\newcommand{\Q}{\mathbb{Q}}
\newcommand{\R}{\mathbb{R}}
\newcommand{\C}{\mathbb{C}}
\newcommand{\K}{\mathbb{K}}
\renewcommand{\L}{\mathbb{L}}
\newcommand{\U}{\mathbb{U}}
\newcommand{\st}{^*}
\newcommand{\tq}{\middle|}
\newcommand{\inv}{^{-1}}
\newcommand{\E}{\mathbb{E}}
\newcommand{\F}{\mathbb{F}}
\newcommand{\V}{\mathbb{V}}
\renewcommand{\ker}{\text{Ker}}
\newcommand{\im}{\text{Im}}
\newcommand{\card}{\text{Card}}
\newcommand{\Tr}{\text{Tr}}

\theoremstyle{plain}
\newtheorem{thm}{Théorème}
\newtheorem*{ind}{Indication}

\theoremstyle{definition}
\newtheorem{dfn}{Definition}

\theoremstyle{remark}
\newtheorem*{rmq}{Remarque}
\newtheorem*{app}{Application}

\newtcolorbox[auto counter]{ex}[2][]{colback=white, colframe=green!50!white, coltitle=black, fonttitle=\bfseries, title={Exercice~\thetcbcounter~: #2}, after title={\hfill #1}}

\title{TD Calcul différentiel}

\begin{document}

\maketitle

Dans tous les exercices, $E$ et $F$ désignent des espaces vectoriels de dimension finie.

\section{Différentiabilité}

\begin{ex}[Moulin.01]{}
	Pour chacune des fonctions suivantes de $\R^2$ dans $\R$, étudier au point $(0,0)$ sa continuité, sa dérivabilité selon tout vecteur de $\R^2$, puis sa différentiabilité.
	\begin{enumerate}
		\item $\displaystyle f_1:(x,y)\mapsto\begin{cases}
			\displaystyle\frac{x^2y}{x^4+y^2}&\text{si }(x,y)\ne(0,0)\\
			0&\text{sinon}
		\end{cases}$
		\item $\displaystyle f_2:(x,y)\mapsto\begin{cases}
			\displaystyle\frac{\sin\lvert xy\rvert}{\lvert x\rvert+\lvert y\rvert}&\text{si }(x,y)\ne(0,0)\\
			0&\text{sinon}
		\end{cases}$
		\item $\displaystyle f_3:(x,y)\mapsto\begin{cases}
			\displaystyle(x+y)\sqrt{x^2+y^2}\sin\left(\frac{1}{\sqrt{x^2+y^2}}\right)&\text{si }(x,y)\ne(0,0)\\
			0&\text{sinon}
		\end{cases}$
		\item $\displaystyle f_4:(x,y)\mapsto\begin{cases}
			\displaystyle\frac{e^{x^4}-e^{y^4}}{(x^2+y^2)^{3/2}}&\text{si }(x,y)\ne(0,0)\\
			0&\text{sinon}
		\end{cases}$
	\end{enumerate}
\end{ex}

\begin{ex}[Moulin.02]{$\circ$ Application de la règle de la chaîne}
	Soit $f:\R^2\to\R$ une fonction de classe $\mathcal{C}^1$.
	\begin{enumerate}
		\item Justifier la dérivabilité de $g:x\mapsto f(x,f(x,x))$.
		\item Déterminer les dérivées partielles de $g:(x,y)\mapsto f(f(x,y),y)$.
		\item Déterminer les dérivées partielles de $g:(x,y)\mapsto f(y,f(x,y))$.
	\end{enumerate}
\end{ex}
\begin{proof}
	Appliquer la règle de la chaîne...
\end{proof}

\begin{ex}[Moulin.03]{}
	\begin{enumerate}
		\item Déterminer le domaine de définition de $f:(x,y)\mapsto\displaystyle\arccos\left(\frac{1-xy}{\sqrt{1+x^2+y^2+x^2y^2}}\right)$.
		\item Déterminer un ouvert $U$ de $\R^2$ sur lequel $f$ est de classe $\mathcal{C}^1$, et calculer les dérivées partielles $\displaystyle\frac{\partial f}{\partial x}$ et $\displaystyle\frac{\partial f}{\partial y}$ sur celui-ci.
		\item En déduire une expression simplifiée de $f$.
	\end{enumerate}
\end{ex}

\begin{ex}[Moulin.04]{}
	Soit $(E,N)$ un espace vectoriel normé de dimension finie. \\
	Etudier la différentiabilité en $0$ de l'application $N$.
\end{ex}

\begin{ex}[Moulin.05]{}
	Etudier les points de différentiabilité sur $\R^n$ des applications $\lVert\cdot\rVert_\infty$ et $\lVert\cdot\rVert_1$.
\end{ex}

\begin{ex}[Moulin.06]{$\star$}
	Soit $E=\R_+\st\times\R_+\st$. Pour $(x,y)\in E$, on pose : $$f(x,y)=\sum_{n=0}^{+\infty}\frac{(-1)^n}{x+ny}$$
	\begin{itemize}
		\item Montrer que $f$ est bien définie sur $E$ et que la série converge uniformément sur tout ensemble de la forme $[a,+\infty[^2$ avec $a>0$. En déduire que $f$ est continue sur $E$.
		\item Montrer que $f$ est de classe $\mathcal{C}^1$.
	\end{itemize}
\end{ex}
\begin{proof}
	Exercice difficile.
	\begin{enumerate}
		\item Utiliser le théorème des séries alternées.
		\item Dériver par rapport à $x$ et à $y$. Pour cela, établir une convergence uniforme sur tout segment : pour la dérivée par rapport à $x$, pas de problème, on a aisément une convergence normale. Pour la dérivée par rapport à $y$, cela est plus difficile : montrer que la suite décroît asymptotiquement en étudiant la fonction et en montrant une borne qui convient pour tout un segment. Ensuite, utiliser le théorème des séries alternées pour obtenir une convergence normale sur tout segment. Ensuite, montrer que ces dérivées partielles sont continues par rapport au couple $(x,y)$. Pour la première, c'est facile, pour la deuxième, réutiliser ce qui a été fait pour la dérivation.
	\end{enumerate}
\end{proof}

\begin{ex}[Moulin.07]{}
	On considère la série de fonctions définie par : $$u_n(x,y)=\frac{x^n}{1+y^{2n}}$$ et $E\subset\R^2$ son ensemble de convergence.
	\begin{enumerate}
		\item Déterminer $E$ et montrer que c'est un ouvert de $\R^2$.
		\item $\star$ Montrer que la somme $f(x,y)=\displaystyle\sum_{n=0}^{+\infty}u_n(x,y)$ est de classe $\mathcal{C}^1$ sur $E$.
		\item Montrer que si $\lvert x\rvert<1$ et $\lvert y\rvert<1$, alors on a : $$f(x,y)=\frac{1}{2}+\sum_{p=0}^{+\infty}\frac{(-1)^pxy^{2p}}{1-xy^{2p}}$$
		\item $\star\star$ En déduire que $f$ n'a pas d'extremum local sur $E$. On étudier d'abord le cas $\lvert y\rvert=1$, puis on ramènera le cas $\lvert y\rvert>1$ au cas $\lvert y\rvert <1$.
	\end{enumerate}
\end{ex}

\begin{ex}[Moulin.08]{}
	\begin{enumerate}
		\item Déterminer l'ensemble de définition $A$ de $f:(x,y)\mapsto\displaystyle\int_{0}^{\pi}\ln(x+y\cos t)dt$.
		\item L'application $f$ est-elle continue sur $A$ ?
		\item L'application $f$ est-elle de classe $\mathcal{C}^1$ sur l'intérieur de $A$ ?
	\end{enumerate}
\end{ex}

\begin{ex}[Moulin.09]{$\Delta$ Fonctions homogènes}
	On pose $U=E\setminus\{0\}$. Soit $\alpha\in\R$ et $f:U\to F$ de classe $\mathcal{C}^1$. On dit que $f$ est \textbf{homogène de degré} $\alpha$ lorsqu'elle vérifie : $$\forall x\in U,\ \forall t\in\R_+\st,\ f(tx)=t^\alpha f(x)$$
	\begin{enumerate}
		\item Montrer que $f$ est homogène de degré $\alpha$ si, et seulement si, on a : $$\forall x\in U,\ \forall t\in\R_+\st,\ df(x)\cdot x=\alpha f(x)$$
		\item Montrer que si $f\in\mathcal{C}^1(U,F)$ est homogène de degré $\alpha$, alors $df$ est homogène de degré $\alpha-1$.
	\end{enumerate} 
\end{ex}
\begin{proof}
	A connaître, très classique !
	\begin{enumerate}
		\item Dans le sens direct, dériver l'hypothèse par rapport à $t$ à $x$ fixé puis prendre $t=1$. Dans le sens réciproque, appliquer l'hypothèse en $tx$ en considérant $g(t)=f(tx)$ pour obtenir une équation différentielle sur $g$ qui se résout en $g(t)=t^\alpha\times C$. Avec $t=1$, on trouve $C=f(x)$.
		\item Différencier l'hypothèse.
	\end{enumerate}
\end{proof}

\begin{ex}[Moulin.10]{$\Delta$ Différentiation des coefficients du polynôme caractéristique}
	On fixe $n\in\N\st$.
	\begin{enumerate}
		\item Pour $M\in\mathcal{M}_n(\C)$, on note $\det(I_n+XM)=\displaystyle\sum_{k=0}^{+\infty}c_k(M)X^k$. montrer que $c_k$ est de classe $\mathcal{C}^1$ pour tout $k\in\N$.
		\item Dans la suite, on se propose de calculer la différentielle de $c_k$. On fixe $A$ et $B$ dans $\mathcal{M}_n(\C)$. \\
		Montrer que pour tout $t$ voisin de $0$, on a : $$\frac{d}{ds}\left(\det(I_n+t(A+sB))\right)\bigg|_{s=0}=\det(I_n+tA)\sum_{k=1}^{+\infty}(-1)^{k-1}\Tr(A^{k-1}B)t^k$$
		\item En déduire, pour tout $k\in\N$ : $$dc_k(A)\cdot B=\sum_{i=1}^{k}(-1)^{i-1}\Tr(A^{i-1}B)c_{k-i}(A)$$
		\item En appliquant la formule précédente à $A=B$, retrouver les formules de Newton.
	\end{enumerate}
\end{ex}

\begin{ex}[Moulin.11 / X]{$\Delta$ Matrices cycliques}
	On considère : $$\begin{array}{lrcl}
		\varphi:&\mathcal{M}_n(\R)&\longrightarrow&\R^n \\
		&M&\longmapsto&(\Tr(M),\ldots,\Tr(M^n))
	\end{array}$$
	\begin{enumerate}
		\item Montrer que $\varphi$ est de classe $\mathcal{C}^1$ et calculer sa différentielle en tout point.
		\item Montrer que pour tout $M\in\mathcal{M}_n(\R)$, on a $\text{rg}(d\varphi(M))=\deg(\pi_M)$.
		\item En déduire que l'ensemble : $\left\{M\in\mathcal{M}_n(\R)\ : \ \chi_{_M}=\pi_M\right\}$ est un ouvert de $\mathcal{M}_n(\R)$.
	\end{enumerate}
\end{ex}
\begin{proof}
	Le nom de l'exercice vient du fait que le dernier ensemble est aussi l'ensemble des matrices $M$ telles qu'il existe un vecteur $X$ tel que $(X,\ldots,M^{n-1}X)$ soit une base de $\R^n$, qu'on appelle ensemble des matrices cycliques.
	\begin{enumerate}
		\item Pour montrer la classe $\mathcal{C}^1$, utiliser les théorèmes généraux.
		\item Utiliser le gradient.
		\item Caractériser la liberté d'une famille par l'inversibilité d'un matrice.
	\end{enumerate}
\end{proof}

\begin{ex}[Morlot / Nougayrède]{}
	Soit $E$ et $F$ deux espaces vectoriels normés de dimension finie. Soit $(f_n)$ une suite de fonctions de classe $\mathcal{C}^1$ de $E$ dans $F$. On suppose que $\sum f_n$ converge simplement sur $E$ et que $\displaystyle\sum(x\mapsto\lVert df(x)\rVert)$ converge normalement au voisinage de tout point. \\ Montrer que $\displaystyle\sum_{n=0}^{+\infty}f_n$ est de classe $\mathcal{C}^1$ et donner sa différentielle.
\end{ex}
\begin{app}
	Montrer que $\exp:\mathcal{M}_n(\K)\to\mathcal{M}_n(\K)$ est de classe $\mathcal{C}^1$.
\end{app}

\begin{ex}[Moulin.12]{$\star$ Différentiation de l'exponentielle}
	\begin{enumerate}
		\item Soit $\mathcal{A}$ une $\C$-algèbre de dimension finie ainsi que $a$ et $b$ deux éléments de $\mathcal{A}$ qui commutent. Montrer l'identité : $$\sum_{n=0}^{+\infty}\left(\frac{1}{n!}\sum_{k=0}^{n-1}a^kb^{n-1-k}\right)=e^b\sum_{n=0}^{+\infty}\frac{(a-b)^n}{(n+1)!}$$ On pourra commencer par le cas où $a-b$ est inversible.
		\item Montrer que $\exp:\mathcal{M}_n(\C)\mathcal{M}_n(\C)$ est de classe $\mathcal{C}^1$.
		\item Montrer que pour tout $A\in\mathcal{M}_n(\C)$, la différentielle de $\exp$ en $A$ vaut : $$\exp(G_A)\circ\sum_{n=0}^{+\infty}\frac{(-\text{ad}_A)^n}{(n+1)!}$$ où $G_A:M\mapsto AM$ et $\text{ad}_A:M\mapsto AM-MA$ sont des endomorphismes de $\mathcal{M}_n(\C)$.
		\item Soit $A\in\mathcal{M}_n(\C)$. Montrer que $d\ \exp(A)$ est un automorphisme de $\mathcal{M}_n(\C)$ si, et seulement si, le spectre de $A$ ne contient pas deux éléments distincts $z$ et $z'$ tels que $z-z'\in 2i\pi\Z$.
	\end{enumerate}
\end{ex}
\begin{ex}[Moulin.13]{Fonctions holomorphes}
	Soit $U$ un ouvert de $\C$ ainsi que $f$ définie sur $U$. On pose $P=\text{Re}(f)$ et $Q=\text{Im}(f)$.\\
	Soit $z_0\in U$. Montrer l'équivalence entre : 
	\begin{enumerate}[label=(\roman*)]
		\item $z\mapsto\displaystyle\frac{f(z)-f(z_0)}{z-z_0}$ a une limite finie en $z_0$ ;
		\item $f$ est différentiable en $z_0$ et $\displaystyle\begin{cases}
			\displaystyle\frac{\partial P}{\partial x}(z_0)=\frac{\partial Q}{\partial y}(z_0) \\
			\displaystyle\frac{\partial Q}{\partial x}(z_0)=-\frac{\partial P}{\partial y}(z_0)
		\end{cases}$
	\end{enumerate}
	Montrer que sous ces conditions, il existe alors un unique complexe $a$ tel que $df(z_0):h\mapsto ah$.
\end{ex}

\begin{ex}[Moulin.14]{}
	Soit $U$ un ouvert convexe de $\R^2$. On considère une fonction $f:U\to\R$ convexe et différentiable.
	\begin{enumerate}
		\item Montrer que $f$ admet un minimum en tout point où sa différentielle est nulle.
		\item montrer que l'ensemble des points où $df$ s'annule est un convexe de $U$ sur lequel $f$ est constante.
	\end{enumerate}
\end{ex}

\begin{ex}[Moulin.15]{$\star$ Théorème fondamental avec hypothèses affaiblies}
	Soit $U$ un ouvert de $\R^2$ et $f:U\to F$. On suppose que $\displaystyle\frac{\partial f}{\partial x}$ et $\displaystyle\frac{\partial f}{\partial y}$ sont définies en tout point, et que la seconde est continue sur $U$. \\ Montrer que $f$ est différentiable en tout point.
\end{ex}

\begin{ex}[Moulin.16]{Inégalité des accroissements finis}
	Soit $U$ un ouvert de $E$ et $f:U\to F$ de classe $\mathcal{C}^1$. Soit $(a,b)\in U^2$ tel que le segment $[a,b]$ soit inclus dans $U$. Montrer l'inégalité : $$\lVert f(b)-f(a)-df(a)[b-a]\rVert\leqslant\underset{x\in[a,b]}{\sup}\lVert(df(x)-df(a))\rVert_{op}\lVert b-a\rVert$$
\end{ex}

\begin{ex}[Moulin.17]{$\star$ Lemme pour l'inversion locale}
	Soit $\Omega$ un ouvert de $E$ contenant $0$ ainsi que $f:\Omega\to E$ de classe $\mathcal{C}^1$ telle que $f(0)=0$ et $df(0)=\text{Id}_E$.
	\begin{enumerate}
		\item Montrer l'existence d'un réel $r>0$ tel que, pour tout $y\in B_f(0,r/2)$, la fonction $g_y:x\mapsto y+x-f(x)$ soit $1/2$-lipschitzienne de $B=B_f(0,r)$ dans lui-même et en déduire l'injectivité de $f$ sur $B$.
		\item Soit $y\in B_f(0,r/2)$. On pose $x_0=0$ et : $$\forall n\in\N,\ x_{n+1}=g_y(x_n)$$ Montrer la convergence de la série $\sum(x_{n+1}-x_n)$ vers un élément $x\in B$ tel que $f(x)=y$, puis que $f$ induit une bijection de $B$ sur un voisinage de $0$ et que $f\inv$ est continue au voisinage de $0$.
		\item Montrer que $f\inv$ est différentiable en $0$.
	\end{enumerate}
\end{ex}
\begin{proof}
	Exercice difficile. A finir...
	\begin{enumerate}
		\item Par continuité, faire en sorte que $$\lVert df(x)-df(0)\rVert\leqslant\frac{1}{2}$$ Puis tout fonction bien. L'injectivité provient du fait que $g_y$ est alors contractante.
		\item On a un $O\left(\displaystyle\frac{1}{2^n}\right)$ donc une convergence absolue. On obtient, par continuité de $g_y$, le fait que $x=g_y(x)$, ce qui traduit le fait que $y=f(x)$. On a donc le caractère surjectif par cette question, et le caractère injectif par la première question. Pour la continuité, à finir...
		\item A faire...
	\end{enumerate}
\end{proof}

\begin{ex}[Moulin.18]{$\star$ Théorème d'inversion locale}
	Soit $\Omega$ un ouvert de $E$ ainsi que $f:\Omega\to E$ de classe $\mathcal{C}^1$ et $a\in\Omega$ tel que $df(a)\in\mathcal{GL}(E)$.
	\begin{enumerate}
		\item Montrer que $df(x)$ est inversible pour tout $x$ dans un voisinage $\Omega'$ de $a$.
		\item En utilisant l'exercice précédent, montrer que $f$ induit une bijection d'un voisinage $U\subset\Omega'$ de $a$ sur un voisinage $V$ de $b=f(a)$ dont la réciproque est continue sur $V$ et différentiable en $b$, puis différentiable en tout point de $V$.
		\item Montrer que $f\inv$ est de classe $\mathcal{C}^1$.
	\end{enumerate}
\end{ex}
\begin{proof}
	A faire...
	\begin{enumerate}
		\item Utiliser la continuité de la différentielle et le fait que $\mathcal{GL}(E)$ est un ouvert de $\mathcal{L}(E)$.
	\end{enumerate}
\end{proof}

\begin{ex}[Moulin.19]{$\star\star$ Inégalité des accroissements finis inversée}
	On suppose $E$ muni d'une structure euclidienne. Soit $f:E\to\R$ de classe $\mathcal{C}^1$. On suppose que $\lVert\nabla f(x)\rVert\geqslant 1$ pour tout $x$ dans $B(0,1)$. \\
	Montrer qu'il existe $x\in B(0,1)$ tel que $f(x)\geqslant f(0)+1$. On pourra commencer par traiter le cas où $\lVert\nabla f(x)\rVert>1$ pour tout $x$ dans $B(0,1)$.
\end{ex}

\begin{ex}[Moulin.20 / X]{$\star$}
	Soit $f:\R\to\R$ telle que $\forall x\in\R,\ f(x+2\pi)-f(x)\in 2\pi\Z$. Soit $F:\R^2\to\R^2$ telle que : $$\forall (r,q)\in\R_+\times\R,\ F(r\cos q,r\sin q)=(r\cos(f(q)),r\sin(f(q)))$$ A quelle condition (nécessaire et suffisante) la fonction $F$ est-elle différentiable en $(0,0)$ ?
\end{ex}

\begin{ex}[Moulin.21]{}
	Soit $n\in\N\st$. On note $U$ l'ouvert de $\R^n$ constitué des $n$-uplets strictement croissants. Montrer que : $$(x_1,\ldots,x_n)\mapsto\prod_{k=1}^{n}(X-x_k)-X^n$$ est de classe $\mathcal{C}^1$ de $U$ dans $\R_{n-1}[X]$ et qu'en tout point sa différentielle est bijective.
\end{ex}
\begin{proof}
	Appliquer le théorème d'inversion globale (HP certes).
\end{proof}

\section{Fonctions de classe $\mathcal{C}^k$}

\begin{ex}[Moulin.22]{$\circ$}
	Sur $\R^2$, on définit $f(x,y)=\displaystyle4xy\frac{x^2-y^2}{x^2+y^2}$ si $(x,y)\ne(0,0)$ et $f(0,0)=0$.
	\begin{enumerate}
		\item Montrer que $f$ est de classe $\mathcal{C}^1$.
		\item Calculer $\displaystyle\frac{\partial^2 f}{\partial x\partial y}(0,0)$ et $\displaystyle\frac{\partial^2 f}{\partial y\partial x}(0,0)$. Que peut-on en conclure ?
	\end{enumerate}
\end{ex}

\begin{ex}[Moulin.23]{}
	\begin{enumerate}
		\item Montrer que la fonction $f$ définie pour $x\ne 0$ par $f(x)=e^{-1/x^2}$ est prolongeable en une fonction de classe $\mathcal{C}^\infty$ sur $\R$ et que on prolongement, que l'on note également $f$, vérifie $\forall k\in\N,\ f^{(k)}(0)=0$.
		\item $\star$ On définit $g(x,y)=\displaystyle\frac{f(x)f(y)}{f(x)+f(y)}$ si $(x,y)\ne(0,0)$ et $g(0,0)=0$. \\ Montrer que $g$ admet des dérivées partielles à tout ordre sur $\R^2$ mais qu'elle n'est pas continue en $(0,0)$.
	\end{enumerate}
\end{ex}

\begin{ex}[Moulin.24]{Un prolongement $\mathcal{C}^\infty$}
	On note $\Delta=\R^2\setminus\{(x,y)\in\R^2\ : \ x=y\}$. Montrer que la fonction $f$ définie sur $\Delta$ par la formule $$f(x,y)=\frac{\sin(x)-\sin(y)}{x-y}$$ se prolonge en une fonction de classe $\mathcal{C}^\infty$.
\end{ex}
\begin{proof}
	Par trigonométrie, si $(x,y)\in\Delta$, on a : $$g(x,y)=\text{sinc}\left(\frac{x-y}{2}\right)\cos\left(\frac{x+y}{2}\right)$$ Or, le sinus cardinal se prolonge par continuité en $0$, de façon $\mathcal{C}^\infty$ car ce prolongement est alors DSE avec un RCV infini. Comme le cosinus est de classe $\mathcal{C}^\infty$, on obtient le résultat par composition.
\end{proof}

\begin{ex}[Moulin.25]{$\star$ Prolongement $\mathcal{C}^\infty$ : cas général}
	On note $\Delta=\R^2\setminus\{(x,y)\in\R^2\ : \ x=y\}$ et considère une fonction $f:\R\to\C$ de classe $\mathcal{C}^\infty$. Montrer que la fonction $g$ définie sur $\Delta$ par la formule $$g(x,y)=\frac{f(x)-f(y)}{x-y}$$ se prolonge en une fonction de classe $\mathcal{C}^\infty$.
\end{ex}
\begin{proof}
	La ruse employée à l'exercice précédent ne fonctionne plus. Et si on veut le faire à la main, ça va être très moche. Donc, on utilise une classe de fonction. On considère : $$S=\left\{(x,y)\mapsto\int_{0}^{1}\sum_{n\in\N}P_n(t)f^{(n)}(x+t(y-x))\ dt \ \mid \ (P_n)\in\R[X]^{(\N)}\right\}$$ Il est clair que $$g:(x,y)\to\int_{0}^{1}f'(x+t(y-x))\ dt$$ est dedans, et c'est le seul prolongement convenable pour $g$. De plus, en dominant, on voit que toute fonction de $S$ est continue. Enfin, en dominant, on montre que $S$ est stable par dérivation partielle. Cela permet de conclure, conformément au point méthode du cours, que $g$ est de classe $\mathcal{C}^\infty$.
\end{proof}

\begin{ex}[Moulin.26]{$\star$ $\mathcal{C}^p$-difféomorphismes}
	Soit $U$ et $V$ deux ouverts de $\R^n$ et $f:U\to V$ une application de classe $\mathcal{C}^p$ avec $p\geqslant 1$. On suppose que $f$ est bijective et que $f\inv$ est différentiable. Montrer que $f\inv$ est de classe $\mathcal{C}^p$.
\end{ex}
\begin{proof}
	A faire...
\end{proof}

\begin{ex}[Moulin.27]{Indépendance du laplacien vis-à-vis des bases orthonormées}
	On suppose $E$ muni d'une structure euclidienne et de dimension $n$. Soit $f:U\to\R$ une fonction de classe $\mathcal{C}^2$. Montrer que $$\Delta f=\sum_{k=1}^{n}\partial_k\partial_k f$$ ne dépend pas de la base orthonormée dans laquelle on le calcule. On l'appelle le laplacien de $f$.
\end{ex}
\begin{proof}
	Le faire par rapport à la base canonique. Utilise la règle de la chaîne et le faire matriciellement.
\end{proof}

\begin{ex}[Laf]{Laplacien en coordonnées polaires}
	Soit $f:\R^2\to\R$ de classe $\mathcal{C}^2$. On définit son laplacien par $$\Delta f=\frac{\partial^2 f}{\partial x^2}+\frac{\partial^2 f}{\partial y^2}$$ On définit ensuite la fonction $g$ sur $\R_+\times\R$ par : $$g(r,\theta)=f(r\cos(\theta),r\sin(\theta))$$ Montrer la formule : $$\Delta f(r\cos(\theta),r\sin(\theta))=\frac{1}{r}\frac{\partial}{\partial r}\left(r\frac{\partial g}{\partial r}\right)+\frac{1}{r^2}\frac{\partial^2 g}{\partial\theta^2}$$
\end{ex}
\begin{proof}
	Calculer les dérivées partielles premières et secondes de $g$ par rapport à $r$ et $\theta$ en utilisant la règle de la chaîne. Ensuite, sommer les dérivées secondes puis ajuster.
\end{proof}

\begin{ex}[Méchanceté]{$\star$ Laplacien en coordonnées sphériques}
	Soit $f:\R^2\to\R$ de classe $\mathcal{C}^2$. On définit son laplacien par $$\Delta f=\frac{\partial^2 f}{\partial x^2}+\frac{\partial^2 f}{\partial y^2}+\frac{\partial^2 f}{\partial z^2}$$ On définit ensuite la fonction $g$ sur $\R_+\times\R\times\R$ par : $$g(r,\theta,\varphi)=f(r\sin(\theta)\cos(\varphi),r\sin(\theta)\sin(\varphi),r\cos(\theta))$$ Montrer la formule : $$\Delta f=\frac{1}{r^2}\frac{\partial}{\partial r}\left(r^2\frac{\partial g}{\partial r}\right)+\frac{1}{r^2\sin(\theta)}\frac{\partial}{\partial\theta}\left(\sin(\theta)\frac{\partial g}{\partial\theta}\right)+\frac{1}{r^2\sin^2(\theta)}\frac{\partial^2 f}{\partial\varphi^2}$$
\end{ex}
\begin{proof}
	Il s'agit de la version énervée du précédent. Un conseil : le faire calmement. Calculer les dérivées partielles premières et secondes de $g$ par rapport à $r$ et $\theta$ en utilisant la règle de la chaîne. Ensuite, sommer les dérivées secondes puis ajuster.
\end{proof}

\begin{ex}[Moulin.28]{}
	Soit $V$ un espace euclidien et $a\in V$. \\ Donner une condition nécessaire et suffisante pour que $x\mapsto\ln\lVert x-a\rVert$ soit une fonction harmonique, c'est-à-dire qu'elle soit de classe $\mathcal{C}^2$ et de laplacien nul.
\end{ex}

\begin{ex}[Moulin.29]{$\star$ Théorème de Schwarz avec hypothèses affaiblies}
	Soit $U$ un ouvert de $\R^2$ et $f:U\to F$. on suppose que les dérivées partielles premières de $f$ existent partout et sont partiellement continues, et que $\displaystyle\frac{\partial^2 f}{\partial x\partial y}$ est définie sur tout $U$ et continue. Montrer alors que $\displaystyle\frac{\partial^2f}{\partial y\partial x}$ est définie sur tout $U$ et que $\displaystyle\frac{\partial^2f}{\partial y\partial x}=\frac{\partial^2 f}{\partial x\partial y}$. On pourra remarquer que : $$f(x,y)=f(x,b)+(y-b)\int_{0}^{1}\frac{\partial f}{\partial y}(x,(1-t)b+ty)\ dt$$
\end{ex}

\begin{ex}[Moulin.30]{$\star$ Théorème de Schwarz avec hypothèses inhabituelles}
	Soit $U$ un ouvert de $\R^2$ et $f:U\to F$. On suppose que les dérivées partielles secondes $\displaystyle\frac{\partial^2f}{\partial x^2}$, $\displaystyle\frac{\partial^2f}{\partial x\partial y}$ et $\displaystyle\frac{\partial^2f}{\partial y^2}$ existent en tout point et que la dernière est continue en tout point. On suppose que $(0,0)\in U$.
	\begin{enumerate}
		\item Montrer que, pour $(x,y)$ voisin de $(0,0)$ : $$f(x,y)=f(0,0)+df(0,0)[x,y]+\frac{x^2}{2}\frac{\partial^2 f}{\partial x^2}(0,0)+xy\frac{\partial^2}{\partial x\partial t}(0,0)+\frac{y^2}{2}\frac{\partial^2 f}{\partial y^2}(0,0)+o(x^2+y^2)$$ On pourra appliquer la formule de Taylor avec reste intégral à la fonction $y\mapsto f(x,y)$.
		\item En déduire que si les quatre dérivées partielles d'ordre $2$ de $f$ sont définies sur tout $U$ et si les dérivées partielles $\displaystyle\frac{\partial^2f}{\partial x^2}$ et $\displaystyle\frac{\partial^2f}{\partial x^2}$ sont continues, alors : $$\displaystyle\frac{\partial^2f}{\partial y\partial x}=\displaystyle\frac{\partial^2f}{\partial x\partial y}$$
	\end{enumerate}
\end{ex}

\begin{ex}[Moulin.31]{$\star$ Théorème de Poincaré}
	Soit $V$ un ouvert étoilé par rapport à $a$ d'un espace euclidien et $f:V\to V$ une fonction de classe $\mathcal{C}^1$ telle que $df(a)$ soit un endomorphisme symétrique pour tout $a$ (on dit aussi que $f$ est de \textbf{rotationnel} nul). Montrer que $f$ est le gradient d'une fonction $F:V\to\R$ de classe $\mathcal{C}^1$. On pourra remarquer que si $f=\nabla F$, alors pour tout $x\in V$ : $$F(x)-F(a)=\int_{0}^{1}\left(f(a+t(x-a))\mid x-a\right)\ dt$$
\end{ex}
\begin{proof}
	Le faire en passant par les expressions du produit scalaire dans la base. Montrer que les dérivées partielles sont les bonnes. Pour cela, dériver sous le signe intégral composante par composante, utiliser l'hypothèse du rotationnel nul puis observer qu'on est en présence d'une règle de la chaîne par rapport à $t$ à l'envers. Remontrer cette règle de la chaîne et faire une intégration par parties pour conclure.
\end{proof}

\begin{ex}[Nougayrède / X]{$\star$ Jacobienne antisymétrique}
	Soit $f:\R^n\to\R^n$ de classe $\mathcal{C}^2$. Montrer que $J_f(x)$ est antisymétrique pour tout $x\in\R^n$ si, et seulement s'il existe $A\in\mathcal{A}_n(\R)$ et $b\in\R^n$ tels que : $$\forall x\in\R^n,\ f(x)=Ax+b$$
\end{ex}

\begin{ex}[Nougayrède / Révisions]{$\star$ Lemme de Morse}
	\begin{enumerate}
		\item SOit $A\in\mathcal{S}_n^{++}(\R)$. On pose $F=A\inv\mathcal{S}_n(\R)$ et : $$\begin{array}{lrcl}
			\varphi:&\mathcal{GL}_n(\R)\cap F&\longrightarrow&\mathcal{GL}_n(\R)\cap\mathcal{S}_n(\R)\\
			&M&\longmapsto&M^TAM
		\end{array}$$ Montrer que $\varphi$ est différentiable en $I_n$ et que $d\varphi(I_n)$ est un isomorphisme de $F$ sur $\mathcal{S}_n(\R)$.
		\item $\star\star$ Montrer qu'il existe $r>0$ tel que $B(A,r)\subset\mathcal{GL}_n(\R)$ et une fonction $\psi$ de classe $\mathcal{C}^1$ sur $$\mathcal{U}=B(A,r)\cap\mathcal{S}_n(\R)$$ telle que $$\psi(A)=I_n\ \ \text{et}\ \ \forall N\in \mathcal{U},\ \begin{cases}
			\psi(N)\in\mathcal{GL}_n(\R)\cap F\\
			N=\psi(N)^TA\psi(N)
		\end{cases}$$
		\item Soit $f\in\mathcal{C}^2(\R^n,\R)$ telle que $f(0)=0$, $\nabla f(0)=0$ et $H_f(0)\in\mathcal{S}_n^{++}(\R)$. \\
		Montrer qu'il existe $V\subset\R^n$ un voisinage de $0$ et $g\in\mathcal{C}^1(V,\R^n)$ vérifiant : $$\forall x\in V,\ f(x)=\lVert g(x)\rVert^2$$ Indication : on pourra utiliser une formule de Taylor.
	\end{enumerate}
\end{ex}

\section{\textit{Extrema}}

\begin{ex}[Moulin.32]{}
	Etudier les \textit{extrema} locaux des fonctions suivantes sur leur ensemble de définition : $$(x,y)\mapsto x^3-3x+xy^2 \ \ \text{et}\ \ (x,y)\mapsto\frac{1}{1-x}+\frac{1}{1-y}+\frac{1}{x+y}$$
\end{ex}

\begin{ex}[Moulin.33]{}
	Etudier les \textit{extrema} des fonctions suivantes :
	\begin{enumerate}
		\item $(x,y)\mapsto\sqrt{x^2+(y-1)^2}+\sqrt{(x-1)^2+y^2}$ ;
		\item $(x,y)\mapsto -2(x-y)^4+x^4+y^4$ ;
		\item $(x,y)\mapsto x^3+y^3-6(x^2-y^2)$ ;
		\item $(x,y)\mapsto \sin x+\sin y+\cos(x+y)$.
	\end{enumerate}
\end{ex}

\begin{ex}[Moulin.34]{Principe du maximum et fonction harmonique}
	Soit $U$ un ouvert non vide borné de $\R^2$ et $f$ une fonction définie sur l'adhérence $U$, continue sur l'adhérence de $U$ et de classe $\mathcal{C}^2$ sur $U$ à valeurs dans $\R$.
	\begin{enumerate}
		\item Montrer que si $\Delta f$ ne prend que des valeurs strictement positives sur $U$, alors $f$ n'a pas de maximum local sur $U$.
		\item $\star$ On suppose que $f$ est \textbf{harmonique} sur $U$, c'est-à-dire que $\Delta f=0$. Montrer que $f$ admet un maximum et un minimum sur l'adhérence de $U$, et que ceux-ci sont atteints sur la frontière de $U$ (il peuvent être atteints ailleurs néanmoins).
	\end{enumerate}
\end{ex}
\begin{proof}
	La deuxième question est difficile. Bel exercice.
	\begin{enumerate}
		\item Utiliser la matrice hessienne qui ne peut pas être négative.
		\item Pour l'existence, c'est une affaire de compacité. Remarquer qu'il suffit de prouver le résultat pour le maximum, quitte à changer $f$ et $-f$. Considérer $$g_n:x\mapsto f(x)+\frac{1}{n+1}\lVert x\rVert^2$$ Lui appliquer la question précédente pour obtenir un maximum sur la frontière. La frontière est compacte, donc on extrait de cette suite de \textit{maxima} une limite. En passant à la limite, un maximum est atteint en ce point.
	\end{enumerate}
\end{proof}

\begin{ex}[Moulin.35]{}
	\textit{Extrema} de $(x,y)\mapsto xe^y+ye^x$.
\end{ex}

\begin{ex}[Moulin.36]{$\circ$}
	Déterminer les \textit{extrema} globaux de la fonction $f:\R^2\to\R$ définie par : $$\forall (x,y)\in\R^2,\ f(x,y)=x^4+y^4-x(x^2-y^2)$$
\end{ex}

\begin{ex}[Moulin.37]{}
	Pour $a\in\R$, on définit $f_a:\R^2\to\R$ par : $$\forall (x,y)\in\R^2,\ f_a(x,y)=(2y+(1+a)x)e^{\displaystyle-\frac{x^2+xy+(1+a)y^2}{2}}$$ Discuter, selon les valeurs de $a$, des \textit{extrema} globaux de $f_a$.
\end{ex}

\begin{ex}[Moulin.38]{}
	Pour $a\in\R$, on définit $g_a:\R^2\to\R$ par : $$\forall (x,y)\in\R^2,\ g_a(x,y)=(2y+x)e^{\displaystyle-\frac{x^2+xy+(1+a)y^2}{2}}$$ Discuter, selon les valeurs de $a$, des \textit{extrema} globaux de $g_a$.
\end{ex}

\begin{ex}[Moulin.39]{}
	Déterminer les \textit{extrema} de la fonction $f:x\mapsto\displaystyle\sum_{i=1}^{n}x_i\ln(x_i)$ sur le simplexe standard (ouvert) $$\Sigma=\left\{(x_1,\ldots,x_n)\in(\R_+\st)^n\ : \ \sum_{i=1}^{n}x_i=1\right\}$$
\end{ex}

\begin{ex}[Moulin.40]{}
	Soit $A\in\mathcal{S}_n(\R)$ non nulle, $U\in\R^n$ et $a\in\R$. Pour $X\in\R^n$, on pose : $$f(X)=X^TAX-2U^TX+a$$
	\begin{enumerate}
		\item Montrer que $f$ est de classe $\mathcal{C}^1$ et calculer son gradient.
		\item On suppose $A$ inversible. Montrer que $f$ admet un seul point critique $X_0$ et calculer $f(X_0)$ en fonction de $A$, $U$ et $a$.
		\item Montrer que si $A$ est définie positive, alors $X_0$ est un minimum global de $f$.
		\item $\star$ Etudier les \textit{extrema} de $f$ lorsque $A$ n'est pas inversible.
	\end{enumerate}
\end{ex}

\begin{ex}[Moulin.41]{}
	Soit $(ABC)$ un triangle dans $\R^2$ affine euclidien. Déterminer le maximum du produit des distances d'un point $M$ intérieur au triangle aux côtés du triangle.
\end{ex}

\begin{ex}[Moulin.42]{}
	Quels sont les triangles de périmètre maximum inscrit dans un cercle ?
\end{ex}

\begin{ex}[Moulin.42]{Fonctions strictement monotones}
	Soit $E$ un espace vectoriel euclidien. Une fonction $f\in\mathcal{C}^1(E,E)$ est dite strictement monotone lorsqu'il existe $k>0$ tel que : $$\forall (x,y)\in E^2,\ \left(f(x)-f(y)\mid x-y\right)\geqslant k\lVert x-y\rVert^2$$
	\begin{enumerate}
		\item Quelles sont les fonctions $\mathcal{C}^1$ strictement monotone de $\R$ dans $\R$ ? Déterminer leur image.
		\item Soit $f\in\mathcal{C}^1(E,E)$. Montrer que $f$ est strictement monotone si, et seulement s'il existe $k>0$ tel que $$\forall (x,h)\in E^2,\ \left(df_x(h)\mid h\right)\geqslant k\lVert h\rVert^2$$
		\item Montrer que si $f\in\mathcal{C}^1(E,E)$ est strictement monotone, alors sa différentielle est inversible en tout point.
		\item $\star$ Montrer qu'une fonction de classe $\mathcal{C}^2$ de $E$ dans $\R$ dont le gradient est strictement monotone admet un unique minimum.
	\end{enumerate}
\end{ex}

\section{\textit{Extrema} liés}

\begin{ex}[Moulin.44]{}
	Soit $(\alpha_1,\ldots,\alpha_n)\in\R_+^n$ tel que $\displaystyle\sum_{i=1}^{n}\alpha_i=1$. \\
	Etudier les \textit{extrema} de $f:x\mapsto\displaystyle\prod_{i=1}^{n}x_i^{\alpha_i}$ sur $K=\displaystyle\left\{(x_1,\ldots,x_n)\in\R_+^n\ : \ \sum_{i=1}^{n}\alpha_ix_i=1\right\}$. \\
	En déduire : $$\forall (x_1,\ldots,x_n)\in\R_+^n,\ \prod_{i=1}^{n}x_i^{\alpha_i}\leqslant\sum_{i=1}^{n}\alpha_ix_i$$
\end{ex}

\begin{ex}[Moulin.45]{}
	Soit $(a,b,c)\in\R^3$ tel que $a>b>c>0$ et : $$S=\left\{(x,y,z)\in\R^3\\ : \ p(x,y,z)=1\right\}\ \ \text{où}\ \ p(x,y,z)=\frac{x^4}{a^4}+\frac{y^4}{b^4}+\frac{z^4}{c^4}$$ Déterminer les \textit{extrema} de $f:(x,y,z)\mapsto x^2+y^2+z^2$ sur $S$.
\end{ex}

\begin{ex}[Moulin.46]{}
	Soit $K_n=\left\{(x_1,\ldots,x_n)\in\R_+^n\ : \ x_1+\cdots+x_n=1\right\}$. On note $M_n$ le \textit{maximum} de sur $K_n$, s'il existe, de : $$f_n:(x_1,\ldots,x_n)\mapsto\sum_{1\leqslant i,j\leqslant n}x_ix_j(x_i+x_j)^2$$
	\begin{enumerate}
		\item Montrer que $M_n$ est bien défini.
		\item Calculer $M_2$. 
		\item Calculer $M_n$.
	\end{enumerate}
\end{ex}

\begin{ex}[Moulin.47]{$\star$}
	Etudier les \textit{extrema} globaux de $\lvert \sin\rvert$ sur le disque unité fermé de $\C$.
\end{ex}

\begin{ex}[Moulin.48]{}
	On note $\mathcal{SL}_n(\R)$ l'ensemble des matrices $M\in\mathcal{M}_n(\R)$ telles que $\det(M)=1$. On munit $\mathcal{M}_n(\R)$ de sa structure euclidienne canonique.
	\begin{enumerate}
		\item Montrer que le gradient du déterminant en $M\in\mathcal{SL}_n(\R)$ est $M^{-T}$.
		\item Déterminer les \textit{extrema} de la norme sur $\mathcal{SL}_n(\R)$.
	\end{enumerate}
\end{ex}

\begin{ex}[Moulin.49]{}
	Soit $A\in\mathcal{M}_n(\R)$. On définit : $$\begin{array}{lrcl}
		\varphi_A:&\mathcal{O}_n(\R)&\longrightarrow&\R\\
		&O&\longmapsto&\Tr(AO)
	\end{array}$$
	\begin{enumerate}
		\item Montrer que $\varphi_A$ admet un \textit{maximum}.
		\item $\star$ Soit $O\in\mathcal{O}_n(\R)$ une matrice en laquelle $\varphi_A$ admet un \textit{maximum}. \\ Montrer que $AO$ est une matrice symétrique.
	\end{enumerate}
\end{ex}

\begin{ex}[Moulin.50]{}
	\textit{Extrema} sur $\R_+^2$ de $(x,y)\mapsto\displaystyle\frac{xy}{(1+y)(1+y)(1+xy)}$
\end{ex}

\begin{ex}[Moulin.51]{}
	Soit $E$ un espace vectoriel euclidien et $f:E\to\R$ une fonction de classe $\mathcal{C}^1$ telle que $$\frac{f(x)}{\lVert x\rVert}\xrightarrow[\lVert x\rVert\to+\infty]{}+\infty$$
	\begin{enumerate}
		\item Soit $v\in E$. Démontrer que $g:w\mapsto f(x)-(x\mid v)$ admet un \textit{minimum}.
		\item En déduire que l'application $\nabla f$ est surjective.
	\end{enumerate}
\end{ex}

\begin{ex}[Moulin.51]{$\Delta$}
	Démontrer l'inégalité arithmético-géométrique : $$\forall x\in\R_+^n,\ \left(\prod_{i=1}^{n}x_i\right)^{1/n}\leqslant\frac{1}{n}\sum_{i=1}^{n}x_i$$
\end{ex}

\begin{ex}[Moulin.53]{$\Delta\circ$ Théorème spectral}
	Soit $E$ un espace euclidien de dimension non nulle, $S$ sa sphère unité et $u$ un endomorphisme autoadjoint de $E$. On définit $f_u$ de $E$ dans $\R$ par : $$f_u:x\mapsto (u(x)\mid x)$$
	\begin{enumerate}
		\item Montrer que $f$ est de classe $\mathcal{C}^1$ et déterminer son gradient.
		\item Sans utiliser les résultats du chapitre sur les endomorphismes d'un espace vectoriel euclidien concernant les endomorphismes autoadjoints, montrer que $f_{\mid S}$ admet un \textit{maximum} et que s'il est atteint en $a\in S$, alors $\nabla f(a)$ et $a$ sont colinéaires.
		\item Retrouver ainsi le fait que $u$ admet une valeur propre réelle.
	\end{enumerate}
\end{ex}
\begin{proof}
	Calculer le gradient de $f$. Montrer que la restriction de $f_u$ à $S$ admet un maximum et utiliser le théorème d'optimisation sous contrainte en ce point. Conclure. 
\end{proof}

\begin{ex}{Une autre démonstration variationnelle du théorème spectral}
	Soit $A\in\mathcal{S}_n(\R)$ et $B=\text{Diag}(1,\ldots,n)$. On pose : $$\begin{array}{lrcl}
		\varphi:&\mathcal{M}_n(\R)&\longrightarrow&\R\\
		&M&\longmapsto&\lVert M^TAM-B\rVert^2
	\end{array}$$ où $\lVert\cdot\rVert$ est la norme euclidienne usuelle sur $\mathcal{M}_n(\R)$.
	\begin{enumerate}
		\item Calculer la différentielle de $\varphi$ en tout point.
		\item Montrer que la restriction de $\varphi$ à $\mathcal{O}_n(\R)$ possède un \textit{maximum}. On notera $P$ une matrice de $\mathcal{O}_n(\R)$ telle que $$\underset{\mathcal{O}_n(\R)}{\max}\varphi=\varphi(P)$$
		\item Soit $C\in\mathcal{A}_n(\R)$ et $\psi:t\mapsto \varphi(Pe^{tC})$. Montrer que $$\forall t\in\R,\ \varphi(t)\leqslant\varphi(0)$$ et en déduire que $$\left(C\mid P^TAPB-BP^TAP\right)=0$$
		\item En déduire que $PA^TP$ et $B$ commutent, puis conclure.
	\end{enumerate}
\end{ex}

\section{Équations aux dérivées partielles}

\begin{ex}[Moulin.54]{$\circ$}
	Déterminer les fonctions $f:\R^2\to\R$ de classe $\mathcal{C}^1$ telles que $$\frac{\partial f}{\partial x}=x^2-y^2\ \ \text{et}\ \ \frac{\partial f}{\partial y}=-2xy$$
\end{ex}

\begin{ex}[Moulin.55]{$\circ$}
	Déterminer les fonctions $f:\R^2\to\R$ de classe $\mathcal{C}^1$ telles que $$\frac{\partial f}{\partial x}=x^3-y^3\ \ \text{et}\ \ \frac{\partial f}{\partial y}=-x^2y^2$$
\end{ex}

\begin{ex}[Moulin.56]{}
	Résoudre sur $\R^2$ l'équation aux dérivées partielles : $$5\frac{\partial f}{\partial x}-3\frac{\partial f}{\partial y}=x+y$$
\end{ex}

\begin{ex}[Moulin.57]{}
	Résoudre sur $\R^2$ l'équation aux dérivées partielles : $$x\frac{\partial f}{\partial x}-y\frac{\partial f}{\partial y}=0$$
\end{ex}

\begin{ex}[Moulin.58]{}
	On considère l'équation aux dérivées partielles $\displaystyle x\frac{\partial f}{\partial x}+y\frac{\partial f}{\partial y}=0$. \\
	Quelles sont les solutions de cette équation sur $\R_+\st\times\R$ ? sur $\R^2\setminus(\R_-\times\{0\})$ ? sur $\R^2\setminus\{(0,0)\}$ ? sur $\R^2$ ?
\end{ex}

\begin{ex}[Moulin.59]{}
	Soit $f\in\mathcal{C}^1(\R^2,\R)$ telle que $(\nabla f(x),x)$ soit liée pour tout $x\in\R^2$. Que dire de $f$ ?
\end{ex}

\begin{ex}[Moulin.60]{$\star$}
	Résoudre sur $\R^2$ l'équation aux dérivées partielles : $$\frac{\partial^2 f}{\partial x^2}+2\frac{\partial^2 f}{\partial y\partial x}-3\frac{\partial^2f}{\partial y^2}=0$$
\end{ex}
\begin{proof}
	Chercher un changement de variables linéaire de la forme. En le cherchant, on trouve qu'il suffit de le prendre sous la forme : $$(x,y)=(u+v,3u-v)$$ pour se ramener à l'exemple du cours avec un $$\frac{\partial^2g}{\partial u\partial v}=0$$
\end{proof}

\section{Vecteurs tangents}

\begin{ex}[Moulin.61]{}
	Déterminer, en tout point des parties suivantes de $\R^3$, l'ensemble des vecteurs tangents, et préciser sa structure selon le point considéré.
	\begin{enumerate}
		\item $(S_1)$ : $z=xy+3y^2$ ;
		\item $(S_2)$ : $z^2=xz+yz-xy$ ;
		\item $(S_3)$ : $x^4+y^4=z^4$.
	\end{enumerate}
\end{ex}

\begin{ex}[Moulin.62]{}
	Soit $S$ la surface de $\R^3$ d'équation $z=xy$ et $M_0=(x_0,y_0,z_0)$ un point de $S$.
	\begin{enumerate}
		\item Déterminer le plan tangent à $S$ en $M_0$.
		\item Montrer que l'intersection de ce plan avec $S$ est la réunion de deux droites affines.
	\end{enumerate}
\end{ex}

\begin{ex}[Moulin.63]{$\star$ Espace tangent à $\mathcal{O}_n(\R)$ en $I_n$}
	Déterminer l'ensemble des vecteurs tangents à $\mathcal{O}_n(\R)$ en $I_n$.
\end{ex}
\begin{proof}
	La réponse est $\mathcal{A}_n(\R)$. Soit $A$ un vecteur tangent à $\mathcal{O}_n(\R)$ en $I_n$. On fixe $\varepsilon>0$ et $$\begin{array}{lrcl}
		\gamma :&]-\varepsilon ,\varepsilon[&\longrightarrow \mathcal{O}_n(\R)
	\end{array}$$ dérivable en $0$ telle que $\gamma(0)=I_n$ et $\gamma'(0)=A$. Alors, en considérant $\varphi:t\mapsto\gamma(t)^T$, qui est alors tout autant définie et dérivable en $0$ que $\gamma$, de dérivée $\varphi'(0)=I_n^T=I_n$. Or : $$\forall t\in]-\varepsilon, \varepsilon[,\ \gamma(t)\varphi(t)=I_n$$ Donc, en dérivant en $0$ : $$0=(\gamma\varphi)'(0)=\gamma'(0)\varphi(0)+\gamma(0)\varphi'(0)=A+A^T$$ si bien que $A\in\mathcal{A}_n(\R)$. Réciproquement, soit $A\in\mathcal{A}_n(\R)$. On considère alors : $$\begin{array}{lrcl}
		\gamma:&]-\varepsilon,\varepsilon[&\longrightarrow&\mathcal{O}_n(\R) \\
		&t&\longmapsto&\exp(tA)
	\end{array}$$ En effet, cette application est déjà bien définie puisque la transposée de l'exponentielle est l'exponentielle de la transposée et alors, puisque $A$ est antisymétrique, $\exp(tA)$ est bien dans $\mathcal{O}_n(\R)$ (en fait dans $\mathcal{SO}_n(\R)$ même avec la formule $\det(\exp)=\exp(\Tr)$). Ensuite, $\gamma$ est dérivable de dérivée $t\mapsto A\gamma(t)$ (\textit{cf} chapitre d'équations différentielles). Ainsi, $\gamma$ est en particulier dérivable en $0$ et : $$\begin{cases}
		\gamma(0)=I_n \\
		\gamma'(0)=A
	\end{cases}$$ si bien que $A$ est effectivement un vecteur tangent à $\mathcal{O}_n(\R)$ en $I_n$, ce qui conclut la preuve.
\end{proof}

\begin{ex}[Moulin.64]{}
	Soit $G$ un sous-groupe de $\mathcal{GL}_n(\K)$ avec $\K=\R$ ou $\K=\C$, et $g\in G$.
	\begin{enumerate}
		\item $\circ$ Déterminer, en fonction de l'ensemble $T$ des vecteurs tangents à $G$ en $I_n$, celui des vecteurs tangents à $G$ en $g$.
		\item Montrer que $T$ est un sous-espace vectoriel de $\mathcal{M}_n(\K)$.
	\end{enumerate}
\end{ex}

\begin{ex}[Moulin.65]{}
	On fixe trois entiers strictement positifs $n$, $p$ et $r$ tels que $\min(n,p)>r$. On note $\mathcal{G}_r$ l'ensemble des matrices de rang $r$ dans $\mathcal{M}_{n,p}(\R)$ et $J_r=\displaystyle\begin{pmatrix}
		I_r&0\\
		0&0
	\end{pmatrix}\in\mathcal{M}_{n,p}(\R)$.
	\begin{enumerate}
		\item Montrer que l'ensemble $$\left\{\begin{pmatrix}
			A&C\\
			B&BA\inv C
		\end{pmatrix}\bigg| (A,B,C)\in\mathcal{GL}_r(\R)\times\mathcal{M}_{n-r,r}(\R)\times\mathcal{M}_{r,p-r}(\R)\right\}$$ est un voisinage de $J_r$ dans $\mathcal{G}_r$.
		\item En déduire l'ensemble des vecteurs tangents à $\mathcal{G}_r$ en $J_r$.
		\item Plus généralement, montrer que pour tout $M_0\in\mathcal{G}_r$, l'ensemble des vecteurs tangents à $\mathcal{G}_r$ en $M_0$ est l'ensemble des matrices $N$ telles que : $$\forall X\in\ker M_0,\ NX\in\text{Im}M_0$$
	\end{enumerate}
\end{ex}

\begin{ex}[Moulin.66]{Espace tangent à $\mathcal{SL}_n(\K)$ en $I_n$}
	L'ensemble des vecteurs tangents à $\mathcal{SL}_n(\K)$ en $I_n$ est l'ensemble des matrices de trace nulle.
\end{ex}
\begin{proof}
	On utilise le théorème du cours avec l'application $g:A\mapsto\det(A)-1$ qui est de classe $\mathcal{C}^1$ et dont l'annulation correspond à $\mathcal{SL}_n(\K)$. Sa différentielle en $I_n$ est celle du déterminant et est bien non nulle : $$dg(I):H\mapsto\Tr(H\text{Com}(I)^T)=\Tr(H)$$ Par résultat de cours, l'espace tangent à $\mathcal{SL}_n(\K)$ en $I_n$ est donc le noyau de cette différentielle, \textit{ie} l'ensemble des matrices de trace nulle.
\end{proof}
\begin{rmq}
	L'ensemble des matrices de trace nulle est de dimension $n^2-1$ si on travaille dans $\mathcal{M}_n(\R)$ et $2n^2-2$ si on travaille dans $\mathcal{M}_n(\C)$ (toujours en tant que $\R$-espaces vectoriels dans ce chapitre). Pour le montrer, appliquer le théorème du rang.
\end{rmq}

\begin{ex}[Moulin.67]{$\star$ Espace tangent à $\mathcal{U}_n(\C)$ en $I_n$}
	Déterminer l'ensemble des vecteurs tangents à $\mathcal{U}_n(\C)$ en $I_n$ et en préciser la dimension.
\end{ex}
\begin{proof}
	La réponse est l'ensemble des matrices anti-hermitiennes, \textit{ie} égales à l'opposé de leur transconjuguée. Adapter la preuve donnée pour $\mathcal{O}_n(\R)$.
\end{proof}

\begin{ex}[Moulin.68]{$\star$}
	Soit $S$ la nappe paramétrée par : $$x=\frac{u+v}{1-uv}\ \ \ y=\frac{1-uv}{1+uv}\ \ \ z=\frac{u-v}{1+uv}$$ Déterminer l'ensemble des vecteurs tangents à $S$ au point de paramètre $(u,v)$ et en donner une équation. \\
	On pourra commencer par déterminer une équation cartésienne de $S$.
\end{ex}

\begin{ex}[Moulin.69]{}
	On considère la surface $S$ d'équation $z=x^4-x^3+xy-y^2$.
	\begin{enumerate}
		\item Déterminer les points où $S$ admet un plan tangent parallèle à $(O;\overrightarrow{i},\overrightarrow{j})$.
		\item Déterminer la position de $S$ par rapport à son plan tangent en ces points.
		\item Déterminer l'intersection de $S$ avec le plan $(O;\overrightarrow{i},\overrightarrow{j})$.
	\end{enumerate}
\end{ex}

\end{document}